% !PRINTED 71
\weberpage{71}

In der Darstellung \S.\,20, (1) von \(\Delta\) werden durch die oberen Indices die Zeilen, durch die unteren Indices die Colonnen gekennzeichnet, und demnach können wir diesem Satze auch den folgenden Ausdruck geben:

\gesperrt{II. Eine Determinante ändert sich nicht, wenn die Zeilen zu Colonnen und die Colonnen zu Zeilen gemacht werden.}

Wenn wir in den sämmtlichen Anordnungen \(\mathfrak A,\mathfrak A',\mathfrak A'',\ldots\) der \(n\) Elemente irgend zwei Elemente mit einander vertauschen, so bleibt die Gesammtheit dieser Anordnungen ungeändert, aber es geht jede Anordnung erster Art in eine Anordnung zweiter Art über und umgekehrt. Wenn wir also in den Gliedern \(M,M',M'',\ldots\), aus denen \(\Delta\) zusammengesetzt ist, irgend zwei untere Indices vertauschen, so geht jedes Glied mit positivem Zeichen in ein anderes über, das in \(\Delta\) mit dem negativen Zeichen behaftet war und umgekehrt, also es ändert \(\Delta\) sein Vorzeichen. Daraus folgt mit Hülfe von II. der Satz:

\gesperrt{III. Wenn man in \(\Delta\) zwei untere oder zwei obere Indices mit einander vertauscht, so ändert die Determinante nur ihr Vorzeichen.}

Etwas anders ausgedrückt: \gesperrt{wenn man zwei Zeilen oder zwei Colonnen mit einander vertauscht, so ändert die Determinante nur ihr Vorzeichen}; und daraus allgemeiner:

\gesperrt{IV. Wenn in einer Determinante die Zeilen oder die Colonnen permutirt werden, so ändert sich der absolute Werth nicht, und das Vorzeichen ändert sich nicht oder geht in das entgegengesetzte über, je nachdem die angewandte Permutation zur ersten oder zweiten Art gehört.}

Aus III. erhält man den folgenden Fundamentalsatz:

\gesperrt{V. Wenn in zwei Zeilen oder in zwei Colonnen die an gleicher Stelle stehenden Glieder einander gleich sind (kürzer ausgedrückt: wenn zwei Reihen einander gleich sind), so hat die Determinante den Werth Null.}

Denn die Vertauschung der zwei Reihen ändert nach III. das Zeichen, kann aber andererseits, da beide Reihen identisch
%CONT>>
